In concluding, we briefly describe how this study is useful for two central enforcement tasks: screening industries for potential cases and writing guidelines to educate companies.

While more work needs to be done, the initial evidence presented here supports a valuable role for an LLM in a screening process. An LLM can substitute for human experts in 
performing an initial sifting through of a large corpus of transcripts. However, given the high rate of false positives and the importance of a competition agency avoiding 
the allocation of resources to false leads, human oversight of the LLM's output is critical. Here is a candidate protocol:

\begin{description}
    \item[Step 1:] LLM reads and rates a corpus of transcripts.
    \item[Step 2:] For those transcripts with the highest scores, human experts read the LLM-generated excerpts and identify those with collusive content.
    \item[Step 3:] For a transcript surviving step 2, human experts read the entire transcript and determine if it is worthy of further examination.
    \item[Step 4:] For those deemed worthy of further examination, all transcripts are collected for the company along with competitors' transcripts and read by human experts.
\end{description}

\noindent Based on the assessment of that body of transcripts, it can be determined if an official investigation should be initiated. As the LLM is refined and 
its performance improved, we expect it to play an expanding role in this screening process.

Prior to pursuing cases based on public communications, it would be prudent for a competition agency to release guidelines describing the content that companies 
should avoid in their public communications. While courts have recognized that public communications can result in an unlawful agreement, exclusive reliance on 
public communications to obtain a conviction is relatively unexplored territory. Issuing guidelines lays the legal groundwork for public prosecution or private 
litigation as well as deterring collusive conduct. As described in \textcite{harrington2022collusion}, a firm's public announcement risks facilitating collusion when it refers 
to the conduct of rival firms or the industry large. Our categorization of the content of the LLM-identified transcripts could be the basis for guidelines 
describing the various types of content that firms should avoid in their public statements.
